% GENERATED FILE. Edit canonical lessons, then run npm run generate:books.
% canonical-source-hash: fnv1a64:8c33aa564ba4ac43
% canonical-lessons: RU-C501-nitka, RU-C501-igolka, RU-C501-napyorstok, RU-C501-vykroyka, RU-C501-podshivat

\chapter{Needle and Thread}
\label{ch:ru-needle-and-thread}
\hlchaptermodality{501}

\begin{chapteropening}
\textbf{By the end of this chapter:} \emph{I can talk about sewing and mending clothes.}

Complete the last lesson of chapter 501: I can talk about sewing and mending clothes.
\end{chapteropening}

\section[\texorpdfstring{\ru{нитка} (nítka)}{нитка (nítka)}]{\ru{нитка} (nítka) — a thread}

\label{lesson:RU-C501-nitka}



\begin{quote}
Before the new one: say the Russian for to embroider, then the Russian for lace.
\end{quote}

\subsection*{You'll want to know: \ru{нитка}}
\textbf{\ru{нитка}} (\emph{nítka}) — \textquotedblleft{}a thread\textquotedblright{}.

\begin{quote}
\ru{Нужна} \ru{чёрная} \ru{нитка}, \ru{чтобы} \ru{пришить} \ru{пуговицу}.
\end{quote}

\begin{quote}
\emph{Nuzhná chyórnaya nítka, chtóby prishít púgovitsu.}
\end{quote}

\begin{quote}
I need black thread to sew the button on.
\end{quote}

Feminine; genitive plural \ru{ниток}. A diminutive of \emph{\ru{нить}}.

\begin{cousinweb}
\emph{\ru{нитка}} is a diminutive of \emph{\ru{нить}} 'thread'; EtymDB lists \emph{\ru{нитка}} without a source, and \emph{\ru{нить}} is inherited from Proto-Slavic \emph{nit\ru{ь}} (EtymDB).
\end{cousinweb}

\begin{grammarlens}[title={What you've built}]
One more word for telling, explaining and describing at B1.
\end{grammarlens}

\subsection*{Your turn}
\begin{itemize}
\raggedright
  \item \emph{Say it:} \emph{nítka}
  \item \emph{Say it:} \emph{nítka}, once more
  \item \emph{Say it:} a sentence of your own with \emph{nítka}
  \item \emph{Recall:} say the Russian for to embroider, then the Russian for lace, then say \emph{nítka} again
\end{itemize}

\begin{tcolorbox}[breakable,colback=teal!4,colframe=teal!35!black,title={Before you move on}]
What does \ru{нитка} mean? (\textquotedblleft{}A thread\textquotedblright{}.) Say it once more.
\end{tcolorbox}
\section[\texorpdfstring{\ru{иголка} (igólka)}{иголка (igólka)}]{\ru{иголка} (igólka) — a needle}

\label{lesson:RU-C501-igolka}



\begin{quote}
Before the new one: say the Russian for lace, then the Russian for a thread.
\end{quote}

\subsection*{You'll want to know: \ru{иголка}}
\textbf{\ru{иголка}} (\emph{igólka}) — \textquotedblleft{}a needle\textquotedblright{}.

\begin{quote}
\ru{Я} \ru{укололся} \ru{иголкой}.
\end{quote}

\begin{quote}
\emph{Ya ukolólsya igólkoy.}
\end{quote}

\begin{quote}
I pricked myself with a needle.
\end{quote}

Feminine. Everyday word beside the bookish \emph{\ru{игла}}.

\begin{cousinweb}
\emph{\ru{иголка}} is a diminutive of \emph{\ru{игла}} 'needle'; EtymDB lists \emph{\ru{иголка}} without a source, and \emph{\ru{игла}} is inherited from Proto-Slavic \emph{j\ru{ь}g\ru{ъ}la} (EtymDB).
\end{cousinweb}

\begin{grammarlens}[title={What you've built}]
One more word for telling, explaining and describing at B1.
\end{grammarlens}

\subsection*{Your turn}
\begin{itemize}
\raggedright
  \item \emph{Say it:} \emph{igólka}
  \item \emph{Say it:} \emph{igólka}, once more
  \item \emph{Say it:} a sentence of your own with \emph{igólka}
  \item \emph{Recall:} say the Russian for lace, then the Russian for a thread, then say \emph{igólka} again
\end{itemize}

\begin{tcolorbox}[breakable,colback=teal!4,colframe=teal!35!black,title={Before you move on}]
What does \ru{иголка} mean? (\textquotedblleft{}A needle\textquotedblright{}.) Say it once more.
\end{tcolorbox}
\section[\texorpdfstring{\ru{напёрсток} (napyórstok)}{напёрсток (napyórstok)}]{\ru{напёрсток} (napyórstok) — a thimble}

\label{lesson:RU-C501-napyorstok}



\begin{quote}
Before the new one: say the Russian for a thread, then the Russian for a needle.
\end{quote}

\subsection*{You'll want to know: \ru{напёрсток}}
\textbf{\ru{напёрсток}} (\emph{napyórstok}) — \textquotedblleft{}a thimble\textquotedblright{}.

\begin{quote}
\ru{Без} \ru{напёрстка} \ru{я} \ru{всегда} \ru{колю} \ru{палец}.
\end{quote}

\begin{quote}
\emph{Bez napyórstka ya vsegdá kolyú pálets.}
\end{quote}

\begin{quote}
Without a thimble I always prick my finger.
\end{quote}

Masculine; genitive \ru{напёрстка}.

\begin{cousinweb}
\emph{\ru{напёрсток}} is inherited from Proto-Slavic \emph{nap\ru{ь}rst\ru{ъ}k\ru{ъ}}, derived from \emph{p\ru{ь}rst\ru{ъ}} 'finger' (EtymDB): a cap on the finger. The old word \emph{\ru{перст}} 'finger', which EtymDB marks as a Church Slavonic borrowing, is usually linked to \emph{\ru{перчатки}} 'gloves'.
\end{cousinweb}

\begin{grammarlens}[title={What you've built}]
One more word for telling, explaining and describing at B1.
\end{grammarlens}

\subsection*{Your turn}
\begin{itemize}
\raggedright
  \item \emph{Say it:} \emph{napyórstok}
  \item \emph{Say it:} \emph{napyórstok}, once more
  \item \emph{Say it:} a sentence of your own with \emph{napyórstok}
  \item \emph{Recall:} say the Russian for a thread, then the Russian for a needle, then say \emph{napyórstok} again
\end{itemize}

\begin{tcolorbox}[breakable,colback=teal!4,colframe=teal!35!black,title={Before you move on}]
What does \ru{напёрсток} mean? (\textquotedblleft{}A thimble\textquotedblright{}.) Say it once more.
\end{tcolorbox}
\section[\texorpdfstring{\ru{выкройка} (výkroyka)}{выкройка (výkroyka)}]{\ru{выкройка} (výkroyka) — a sewing pattern}

\label{lesson:RU-C501-vykroyka}



\begin{quote}
Before the new one: say the Russian for a needle, then the Russian for a thimble.
\end{quote}

\subsection*{You'll want to know: \ru{выкройка}}
\textbf{\ru{выкройка}} (\emph{výkroyka}) — \textquotedblleft{}a sewing pattern\textquotedblright{}.

\begin{quote}
\ru{Мама} \ru{сшила} \ru{платье} \ru{по} \ru{старой} \ru{выкройке}.
\end{quote}

\begin{quote}
\emph{Máma sshíla plátye po stároy výkroyke.}
\end{quote}

\begin{quote}
Mum made a dress from an old pattern.
\end{quote}

Feminine. \emph{\ru{Шить} \ru{по} \ru{выкройке}} — to sew from a pattern.

\begin{cousinweb}
\emph{\ru{выкройка}} is built on \emph{\ru{выкроить}} 'to cut out', from \emph{\ru{кроить}} 'to cut cloth'; EtymDB has no entry.
\end{cousinweb}

\begin{grammarlens}[title={What you've built}]
One more word for telling, explaining and describing at B1.
\end{grammarlens}

\subsection*{Your turn}
\begin{itemize}
\raggedright
  \item \emph{Say it:} \emph{výkroyka}
  \item \emph{Say it:} \emph{výkroyka}, once more
  \item \emph{Say it:} a sentence of your own with \emph{výkroyka}
  \item \emph{Recall:} say the Russian for a needle, then the Russian for a thimble, then say \emph{výkroyka} again
\end{itemize}

\begin{tcolorbox}[breakable,colback=teal!4,colframe=teal!35!black,title={Before you move on}]
What does \ru{выкройка} mean? (\textquotedblleft{}A sewing pattern\textquotedblright{}.) Say it once more.
\end{tcolorbox}
\section[\texorpdfstring{\ru{подшивать} (podshivát)}{подшивать (podshivát)}]{\ru{подшивать} (podshivát) — to hem}

\label{lesson:RU-C501-podshivat}



\begin{quote}
Before the new one: say the Russian for a thimble, then the Russian for a sewing pattern.
\end{quote}

\subsection*{You'll want to know: \ru{подшивать}}
\textbf{\ru{подшивать}} (\emph{podshivát}) — \textquotedblleft{}to hem\textquotedblright{}.

\begin{quote}
\ru{Брюки} \ru{слишком} \ru{длинные}, \ru{и} \ru{бабушка} \ru{их} \ru{подшивает}.
\end{quote}

\begin{quote}
\emph{Bryúki slíshkom dlínnyye, i bábushka ikh podshiváyet.}
\end{quote}

\begin{quote}
The trousers are too long, so Grandma is taking them up.
\end{quote}

Imperfective; perfective \ru{подшить} (\ru{подошью}).

\begin{cousinweb}
\emph{\ru{подшивать}} is built on \emph{\ru{под}-} 'under' and \emph{\ru{шить}} 'to sew'; EtymDB lists the verb without a source.
\end{cousinweb}

\begin{grammarlens}[title={What you've built}]
One more word for telling, explaining and describing at B1.
\end{grammarlens}

\subsection*{Your turn}
\begin{itemize}
\raggedright
  \item \emph{Say it:} \emph{podshivát}
  \item \emph{Say it:} \emph{podshivát}, once more
  \item \emph{Say it:} a sentence of your own with \emph{podshivát}
  \item \emph{Recall:} say the Russian for a thimble, then the Russian for a sewing pattern, then say \emph{podshivát} again
\end{itemize}

\begin{tcolorbox}[breakable,colback=teal!4,colframe=teal!35!black,title={Before you move on}]
What does \ru{подшивать} mean? (\textquotedblleft{}To hem\textquotedblright{}.) Say it once more.
\end{tcolorbox}
